\section{Method}
\label{sec:method}

The audit plan---metrics, five hypotheses and decision rules---was committed and
pushed to the study repository (commit \texttt{2b0bc54}, 16:32 on the day of the
audit) before any audit script existed. It is not a registered report: the plan
was written by the same person who built the game and wrote the commits that B5
labels, so B5 and E1 are not blind. Two of the five hypotheses restate facts a
read-only survey had already seen while scoping the study (triangle counts,
texture sizes, the double-sided flag, the root rotation, unit-scale heights);
the plan says so. The scripts and all first results were then committed together
21 minutes later, so the order in which they ran is not separately recorded.
Everything added after the first draft, in response to an internal fact-check and
review, is marked as such and listed in Section~\ref{sec:deviations}.

\paragraph{Terms.} glTF stores a \emph{seam-split} vertex buffer: wherever a
per-vertex attribute such as a UV coordinate is discontinuous, the vertex is
duplicated. \emph{Welding} merges vertices with identical positions and keeps
UVs per corner, as \emph{wedge} attributes~\cite{hoppe1999newqem}. After welding,
a mesh is \emph{closed} if it has no boundary edges and \emph{edge-manifold} if
every edge has exactly two incident triangles. MiB is $2^{20}$ bytes.

\paragraph{B1 Conventions.} A small GLB reader (no glTF library, so the numbers
come from the bytes any importer reads) applies each file's node chain and
records the bounding box, the largest extent, the origin as a fraction of the box
per axis, node transforms, material flags and extensions, and each embedded
image's format and size. \textbf{H1:} every raw asset is scale-normalised (largest
extent 0.9--1.2\,units), carries a non-identity root rotation and is
double-sided. \emph{Added after review:} each GLB was rendered in glTF
coordinates from $+Z$ and $-Z$ to see which side is the front.

\paragraph{B2 Hygiene.} After welding: connected components (triangles sharing a
vertex), boundary and non-manifold edges, non-manifold vertices, degenerate and
duplicate triangles, and self-intersecting triangles (MeshLab~\cite{cignoni2008meshlab}
via PyMeshLab~\cite{muntoni2021pymeshlab}). Before welding: UV charts (seams
separate them) and UV coverage, the share of the unit square covered by UV
triangles rasterised at 2048$^2$. Rasterisation counts edge texels, so coverage
is an upper bound by about two points; chart overlap was not measured.
\textbf{H2:} a majority of the 24 raw assets fail at least one of closed,
edge-manifold, or single component. \emph{Added after review:} the same
measurements on the game's 81 non-generated meshes (a licensed character and three
purchased environment packs).

\paragraph{B3 Cost.} Geometry bytes as stored and texture bytes for three GPU
formats with full mip chains: RGBA8 (4\,B/texel), BC7 (1\,B/texel), BC1
(0.5\,B/texel). \textbf{H3:} texture memory exceeds geometry memory at least
tenfold for every raw asset in every format. \emph{Added after review:} the
textures actually stored in the shipped Windows build, read with UnityPy.

\paragraph{B4 Simplification.} Texture-aware quadric edge
collapse~\cite{garland1997qem,garland1998colorqem,hoppe1999newqem} as implemented
in MeshLab (quality threshold 0.3, boundary preservation off) at 50, 25, 12.5 and
6.25\,\% of the triangles, up to three passes. The plan named a with-UV and a
no-UV condition. Because B2 found the shipped companion torn, the with-UV
condition was run on two inputs: \emph{S}, the seam-split buffer that MeshLab's
glTF loader returns (what our companion pipeline used), and \emph{W}, the same
file after \texttt{meshing\_remove\_duplicate\_vertices}. \emph{N} is welded, no
UVs. We record triangles reached, boundary edges in the output after welding,
and surface distance to the original sampled Metro-style~\cite{cignoni1998metro}
(200,000 samples per direction; the larger of the two directions, as a percentage
of the bounding-box diagonal; we report the mean and the maximum).
\textbf{H4:} with UVs preserved, most assets cannot reach 6.25\,\% within three
passes while the no-UV control can. \emph{Added after review} (12.5 and
6.25\,\%): the same S and W inputs in meshoptimizer~\cite{meshopt}; S with
boundary preservation on; N with the same normal flag as S and W;
and a colour error for S and W, comparing base colour looked up through the
simplified UVs with base colour at the nearest original vertex (60,000 samples;
the same procedure on the unsimplified mesh gives the noise floor).

\paragraph{B5 History.} All 484 non-merge commits up to the final build were
searched with a pre-stated pattern of 25 terms (scale, pivot, ground, clip, cull,
bounds, glb, decimat, \ldots). Each of the 156 matches was labelled
\texttt{ASSET} (root cause is a property of a generated asset),
\texttt{ENGINE} (an engine or tool behaviour any asset would hit),
\texttt{LEVEL} (level geometry, placement, purchased or procedural content) or
\texttt{OTHER}, with ambiguous cases sent to \texttt{OTHER}. Two AI labellers
drafted the labels, each taking half; the AI assistant reviewed every
\texttt{ASSET}, \texttt{ENGINE} and low-confidence label in two rounds, recording
each change with its reason, and the author signed off. For \texttt{ASSET}
commits we note which B1--B3 measurement, if any, would have shown the property
before import. \textbf{H5:} most \texttt{ASSET} fixes concern properties such
counting would reveal.

\paragraph{E1 Grounding (exploratory).} In Blender~5.1 we evaluate the boss's
skinned mesh at every frame of each animation (60\,fps) and take the lowest vertex
as ground truth. Heights are scaled so that the bind pose is 10\,m tall, as in the
game. We score the estimators the game tried: bind-pose bounds, one frame, and
the shipped minimum over ten samples 0.16\,s apart, at every starting phase of
the loop.

\paragraph{C Attribution audit (second plan).} After the first results, a second
plan was committed (\texttt{63b4a34}) before any claim was extracted. 410 of the
484 commits carry the AI assistant's co-author trailer; their messages, the
comments in the asset tools and the design notes explain what the generated
assets are and why problems occurred. The source texts were frozen (174 commits,
215 comment lines, five design notes). Two AI extractors, without access to the
measurements, independently listed every \emph{claim}: a declarative statement
about a property of a generated asset, about what the pipeline does to it, or
about the cause of an asset problem. The merged claims were judged by the
assistant and, independently, by a second AI judge as \texttt{CONFIRMED},
\texttt{CONTRADICTED}, \texttt{PARTLY} (right observation, wrong cause or a
magnitude off by more than $2\times$) or \texttt{UNTESTABLE}, using only this
study's measurements and two measurements listed in the plan in advance (the
origin after UniGLTF's import; what MeshLab's loader does with the root node).
The first judge is the assistant itself. Disagreements were resolved by the
assistant under two rules: a verdict needs a measurement of the asset concerned,
and claims about engine-side readings we did not reproduce are untestable.
Extraction, judging and resolution were committed together, so their order is
not separately recorded in git; the author reviewed the final verdicts. \textbf{H6:} at least 20\,\% of testable claims are contradicted or
partly wrong. \textbf{H7:} most contradicted claims are never revised before the
final build. \textbf{H8} (exploratory): causal claims are wrong more often than
property claims.

\paragraph{D Defect lifecycle.} For every defect the audit identified and every
\texttt{ASSET}/\texttt{ENGINE} fix, we reconstructed from git when it was
introduced, first described (and by whom: play-test, import inspection, the
assistant, an AI review, this audit), fixed, and whether it reached the final
build. \textbf{H9} (descriptive): defects visible on screen were fixed during
development; defects not visible on screen shipped. The interactive session
transcripts of the development period no longer exist (the earliest surviving one
is from September), so effort is measured in commits and elapsed time.
